\paragraph{The three finite invariant factors.}
Retain $A=k[[t]]$, $q=|k|$, odd characteristic, the original even and
odd block spaces $E_A,O_A$, and the representation
$R(g)=(1+\sum_i\varepsilon_i\kappa_i(g))\rho(g)$.
The primal and paired modules are
\[
 M_A=E_A\oplus\bigoplus_i\varepsilon_iO_A,
 \qquad W_A=\eta E_A\oplus\bigoplus_i\varepsilon_i'O_A.
\]
They are integral coordinate direct summands, with their stated
reductions over $A_m=A/(t^m)$.

\begin{lemma}[Counting integral matrix equations]
\label{lem:selmer-finite-invariant-count}
Let a family of integral linear equations on a finite free $A$-module
define an integral kernel of rank $r$. There are nonnegative integers
$a_1,\ldots,a_e$ such that the number of solutions modulo $t^m$ is
\[
                  q^{mr+\sum_{i=1}^e\min(m,a_i)}
                  \leq q^{mr+\sum_i a_i}\qquad(m\geq0).
\]
\end{lemma}
\begin{proof}
The submodule of the finite free dual generated by all equation rows
is finitely generated. Finitely many original rows therefore generate
it and give the same equations at every truncation. Put their matrix
in Smith normal form. A nonzero diagonal entry $t^{a_i}$ allows
$q^{\min(m,a_i)}$ choices, and each of the $r$ free coordinates
allows $q^m$ choices. This counts every congruence solution; no
integral lift of a solution is required.
\end{proof}

\begin{proposition}[Global, distinguished and auxiliary invariants]
\label{prop:selmer-three-invariant-bounds}
Retain the original even-span hypothesis. For the distinguished
subgroup $D$, assume the original condition
\[
 \bigl(\chi(g)\rho(g)Y=Y\rho(g)\quad(g\in D)\bigr)
                    \Longrightarrow Y=0\qquad(Y\in M_n(A)).
\]
Fix an auxiliary element $\sigma$ whose characteristic polynomial
splits into distinct linear factors over the fraction field, with
roots in $A$. There are constants $C_G,C_0,C_\sigma$, independent
of $m$, such that
\[
\begin{aligned}
 \#\{Z\in M_{A_m}:Z^T=-Z,\ [R(g),Z]=0\ (g\in G)\}
      &\leq q^{C_G},\\
 \#\{Z\in W_{A_m}:\chi(g)R(g)ZR(g)^{-1}=Z\ (g\in D)\}
      &\leq q^{C_0},\\
 \#\{Z\in W_{A_m}:[R(\sigma),Z]=0\}
      &\leq q^{mn+C_\sigma}.
\end{aligned}
\]
The last integral invariant module has rank exactly $n$.
\end{proposition}
\begin{proof}
A matrix commuting with $\rho(G)$ commutes with every diagonal
matrix unit: the span hypothesis supplies a nonzero power of $t$
times that unit, which can be cancelled. It is therefore diagonal.
A skew such matrix is zero, since 2 is invertible. For a skew
$Z=Z_0+\sum_i\varepsilon_iZ_i$ commuting with $R(G)$, the constant
coefficient first gives $Z_0=0$, and the $\varepsilon_i$ coefficients
then give $Z_i=0$.

For the distinguished action, write
$Z=\eta Z_0+\sum_i\varepsilon_i'Z_i$. The
$\varepsilon_i'$ coefficients satisfy the displayed distinguished
condition, hence vanish. The remaining $\eta$ coefficient satisfies
the same condition, hence vanishes too.

Put $P=\rho(\sigma)$ and $K=k((t))$. Its simple-spectrum
centralizer over $K$ is commutative of dimension $n$. Each odd
coefficient $Z_i$ of a commuting paired matrix commutes with $P$.
So does the block projector $Q=\operatorname{diag}(1_a,0_{n-a})$.
Thus $Z_i$ commutes with $Q$ and, being off-diagonal, is zero.
The remaining $Z_0$ lies in the centralizer of $P$. Conversely
every member of that centralizer commutes with $Q$, is even, and
gives an invariant $\eta Z_0$. The integral rank is consequently
$n$. Apply Lemma~\ref{lem:selmer-finite-invariant-count} to the three
original systems of equations \cite{alpo-invariant-counts}.
\end{proof}

The constants retain the finite torsion defects of these systems;
zero integral invariants do not imply zero invariants at each
truncation. Eigenvalue differences need not be units. All constants
are for the fixed arithmetic point and chosen elements, without
asserting uniformity over moving points.
